% =====================================================================
%  ABSTRAK — ditulis dalam satu paragraf, memuat masalah, metode,
%  dan kesimpulan. Angka hasil diisi setelah eksperimen selesai.
% =====================================================================

Pencatatan berat barang di gudang dan pabrik masih dilakukan secara manual,
sehingga kesalahan catat tidak terdeteksi dan tidak tersedia jejak audit.
Pengenalan karakter optis untuk tampilan tujuh segmen telah berkembang, tetapi
seluruh penelitian yang ada berhenti pada pelaporan akurasi dan tidak
melaporkan bacaan salah yang lolos tanpa peringatan. Penelitian ini merancang
sistem pembacaan tampilan tujuh segmen dengan empat lapisan verifikasi
kontekstual, yaitu verifikasi struktural, fisik, temporal, dan statistik, di
atas keluaran PP-OCRv6 yang dijalankan tanpa koneksi jaringan. Data latih
dibangkitkan secara sintetis dengan kondisi terkontrol, sedangkan data uji
berupa citra timbangan nyata yang dikunci sebelum eksperimen. Pengujian dilakukan
berjenjang melalui enam eksperimen untuk memisahkan kontribusi setiap lapisan.
Keberhasilan diukur dengan laju kesalahan senyap, \textit{coverage},
\textit{selective accuracy}, dan luas daerah di bawah kurva
\textit{risk--coverage}. Hasil penelitian menunjukkan bahwa sistem menurunkan
laju kesalahan senyap dari \isian{angka E1} menjadi \isian{angka E5} pada
\textit{coverage} sebesar \isian{angka}, dengan waktu pemrosesan
\isian{angka} milidetik per transaksi. Dengan demikian, penelitian ini
memberikan kerangka evaluasi keandalan yang belum tersedia pada literatur
pembacaan tampilan tujuh segmen.

\vspace{1\baselineskip}
\noindent\textbf{Kata kunci:} tampilan tujuh segmen, kesalahan senyap, verifikasi
kontekstual, PP-OCRv6, \textit{selective prediction}
